\documentclass[11pt,a4paper]{article}
\usepackage[T1]{fontenc}
\usepackage{lmodern}
\usepackage[margin=28mm]{geometry}
\usepackage{amsmath,amssymb,amsthm,mathtools,microtype}
\usepackage{enumitem}
\usepackage[numbers,sort&compress]{natbib}
\usepackage{xcolor}
\usepackage[colorlinks=true,linkcolor=blue!40!black,citecolor=blue!40!black,urlcolor=blue!40!black]{hyperref}
\numberwithin{equation}{section}
\newtheorem{theorem}{Theorem}[section]
\newtheorem{lemma}[theorem]{Lemma}
\newtheorem{proposition}[theorem]{Proposition}
\newtheorem{corollary}[theorem]{Corollary}
\theoremstyle{definition}
\newtheorem{definition}[theorem]{Definition}
\theoremstyle{remark}
\newtheorem{remark}[theorem]{Remark}
\newcommand{\norm}[1]{\left\lVert#1\right\rVert}
\newcommand{\dist}{\operatorname{dist}}
\newcommand{\diam}{\operatorname{diam}}
\newcommand{\tr}{\operatorname{tr}}
\newcommand{\Comet}{\mathsf{C}}
\newcommand{\Kite}{\mathsf{K}}
\newcommand{\LL}{\mathcal L_n}
\setlist[enumerate]{itemsep=3pt,topsep=5pt}
\hypersetup{pdftitle={Eventual minimizers of the adjacency spectral gap in trees and connected graphs},pdfauthor={},pdfsubject={Symmetric double comets and double kites}}
\title{\bfseries Eventual minimizers of the adjacency spectral gap\\in trees and connected graphs}
\author{}
\date{}
\begin{document}
\maketitle
\begin{abstract}
The adjacency spectral gap of a connected graph is the difference between
its two largest adjacency eigenvalues. We prove that, for every sufficiently
large order, every tree minimizing this gap is a symmetric double comet,
and every connected graph minimizing it is a symmetric double kite.
These results establish the eventual forms of the conjectures of
Jovovi\'c, Koledin and Stani\'c for trees and of Stani\'c for connected
graphs. The number of leaves at each terminal vertex of a minimizing
tree, and the order of each terminal clique of a minimizing graph, are
asymptotic to $n/(2\log n)$.
The proofs combine resistance bounds with comparisons against explicit
families. For trees, nodal components identify two dominant vertices, and
a Schur complement forces their degrees to agree before an exact path
comparison eliminates all other branches. For graphs, a centered second
eigenvector gives a constant vertex budget, from which two dense cores are
extracted. A comparison at fixed order removes exceptional vertices and
multiple attachments; a uniform completion argument forces both cores to
be cliques. All auxiliary estimates needed for these reductions are proved.
\end{abstract}
\noindent\textbf{Keywords.} Adjacency spectral gap; extremal graphs; trees;
double comets; double kites; Perron vector; Schur complement.\\
\textbf{Mathematics Subject Classification (2020).} 05C50, 05C05, 15A18.

\section{Introduction}

\subsection{Notation and the extremal problems}
All graphs are finite, simple and undirected. The order of a graph $G$ is
$n=|V(G)|$, and its adjacency matrix is $A(G)$. We write its eigenvalues as
\[
 \lambda_1(G)\ge\lambda_2(G)\ge\cdots\ge\lambda_n(G).
\]
For a connected graph of order at least two, its \emph{adjacency spectral
gap} is
\[
 \delta(G)=\lambda_1(G)-\lambda_2(G)>0.
\]
Whenever a graph is fixed, we suppress $G$ in its eigenvalues. Throughout,
$\mathbf{x}$ and $\mathbf{y}$ denote orthonormal real eigenvectors for
$\lambda_1$ and $\lambda_2$, respectively, with $\mathbf{x}>0$. Their
coordinates at a vertex $u$ are $x_u$ and $y_u$. Auxiliary vectors receive
different letters. For any graph $H$, $\rho(H)=\lambda_1(H)$ denotes its
spectral radius. Matrix norms are Euclidean operator norms, and vector
norms without a subscript are Euclidean norms. The notation $M\preceq N$
means that $N-M$ is positive semidefinite. All logarithms are natural.

Let $\mathcal T_n$ and $\mathcal G_n$ be the sets of trees and connected
graphs of order $n$, and put
\[
 \tau_n=\min_{T\in\mathcal T_n}\delta(T),\qquad
 \eta_n=\min_{G\in\mathcal G_n}\delta(G).
\]
The minima exist because the sets of isomorphism classes are finite.
Write $P_j$ for a path on $j$ vertices and $K_j$ for a complete graph on
$j$ vertices. We use \emph{edge length} for every joining path:
\begin{itemize}
\item $\Comet(k,L)$ is the tree obtained from a path of $L\ge1$ edges by
attaching $k\ge0$ new leaves to each endpoint. It is a symmetric double
comet and has $2k+L+1$ vertices.
\item For $k\ge1$, $\Kite(k,L)$ consists of two disjoint copies of $K_k$, with one
selected vertex in each joined by a path of $L\ge1$ edges whose internal
vertices are new. It is a symmetric double kite and has $2k+L-1$ vertices.
\end{itemize}
Thus $\Comet(k,L)$ is $C(k,L+1)$ in the notation of
\citet{JKS2018}, whereas $\Kite(k,L)$ is $DK(k,L-1)$ in the notation of
\citet{Stanic2013}.

\subsection{Previous work}
Stani\'c formulated the conjecture that some double kite attains $\eta_n$
at every order \citep[Conjecture~2.1]{Stanic2013}. That paper verified the
conjecture by exhaustive computation through order ten, reported
counterexample searches through order fifteen, and derived equations for
the first two eigenvalues of double kites, together with spectral-gap
bounds and structural restrictions on possible minimizers. In particular,
its double-kite calculations exhibit the role of the long joining path.

Jovovi\'c, Koledin and Stani\'c proposed the corresponding conjecture that
some symmetric double comet attains $\tau_n$
\citep[Conjecture~2.1]{JKS2018}. They verified it through order twenty.
They also proved that, for a fixed internal path, moving one terminal
leaf from the larger end of a generalized double comet to the smaller
end decreases its gap whenever the two leaf counts differ by at least
two \citep[Theorem~2.3]{JKS2018}. Their further results include eigenvalue
equations and structural conditions on extremal trees. The two
conjectures are included in the survey of Liu and Ning
\citep{LiuNing2023}.

Several related extremal problems concern $\lambda_2$ itself.
Kumar, Mohar, Pragada and Zhan determine the trees maximizing the second
adjacency eigenvalue at prescribed order and diameter, and develop
spectral-center decompositions for the associated extremal problems
\citep{KMPZDiameter}. The same authors study the maximum of nonnegative
convex combinations of the first two eigenvalues over trees of fixed
order \citep{KMPZConvex}. Wang and Guo investigate
extremal second eigenvalues in classes of graphs at prescribed order
\citep{WangGuo2026}. These problems provide related structural settings;
the objective here is the difference $\lambda_1-\lambda_2$, and all
comparisons needed for that objective are established below.

\subsection{Main results and proof strategy}
Our first theorem identifies every minimizing tree at sufficiently large
orders.
\begin{theorem}\label{thm:trees}
There is an integer $N_{\mathcal T}$ such that, for every $n\ge N_{\mathcal T}$,
every $T\in\mathcal T_n$ with $\delta(T)=\tau_n$ is isomorphic to
$\Comet(k,L)$ for some $2k+L+1=n$. Along every sequence of such minimizers,
\[
 k\sim\frac{n}{2\log n},\qquad
 L=n-(1+o(1))\frac{n}{\log n}.
\]
\end{theorem}
The second theorem gives the corresponding conclusion for connected graphs.
\begin{theorem}\label{thm:graphs}
There is an integer $N_{\mathcal G}$ such that, for every $n\ge N_{\mathcal G}$,
every $G\in\mathcal G_n$ with $\delta(G)=\eta_n$ is isomorphic to
$\Kite(k,L)$ for some $2k+L-1=n$. Along every sequence of such minimizers,
\[
 k\sim\frac{n}{2\log n},\qquad
 L=n-(1+o(1))\frac{n}{\log n}.
\]
\end{theorem}
Neither theorem asserts uniqueness of the minimizing integer $k$. The
conclusions concern all sufficiently large orders; the conjectures at
every order are not consequences of these statements. No explicit
threshold is needed in the proofs.

Both arguments begin by comparing a resistance lower bound with a
symmetric comparison family. For trees, this determines the diameter and
spectral scales. Two sign components of $\mathbf y$ then produce vertices
whose degrees exhaust the branching budget to leading order. Deleting
these vertices leaves a matrix of operator norm $o(\lambda_1)$. A two-dimensional
Schur complement places the two integer degrees in a common interval of
length $o(1)$, forcing equality. An exact path-resolvent estimate then
shows that any remaining branch loses to a double comet of the same order.

For general graphs we center $\mathbf y$ by subtracting a multiple of
$\mathbf x$. The centered extrema and an exact path-resistance identity
leave at most two vertices of slack in a neighborhood count. All large
positive and negative coordinates form two dense terminal cores, with a
bounded residual attachment set. A finite-rank resolvent estimate
quantifies the spectral benefit of these attachments, while path
transmission quantifies their effect on the gap. The comparison at fixed
order removes every extra vertex, every unequal core order, and every
multiple attachment. Finally, a uniform rooted-core comparison shows
that any missing clique edge strictly increases the gap.

We use the Perron--Frobenius theorem, the variational principle, eigenvalue
interlacing and elementary symmetric-matrix inequalities in their standard
forms; see \citet{BrouwerHaemers2012}. The path identities, resistance
estimates, and Schur-complement comparisons specific to this problem are
proved in full. Asymptotic statements are always made along an arbitrary
sequence of minimizers whose orders tend to infinity. A contradiction to
an arbitrary sequence of exceptions therefore establishes one threshold
valid for every minimizer.

\section{Resistance, path resolvents and two-vertex elimination}
\label{sec:preliminaries}

\subsection{The ground-state identity}
Put $\pi_u=x_u^2$, so that $\sum_u\pi_u=1$. For a real function $f$ on
$V(G)$, set $z_u=x_uf_u$. The Perron equations give
\begin{equation}\label{pre:energy}
 \mathbf z^T(\lambda_1I-A)\mathbf z
 =\sum_{uv\in E(G)}x_ux_v(f_u-f_v)^2.
\end{equation}
Indeed, expanding the right-hand side and collecting the coefficient of
$f_u^2$ gives $x_u\sum_{v\sim u}x_v=\lambda_1x_u^2$.
The variational principle consequently yields
\begin{equation}\label{pre:variational}
 \delta(G)=\min_{\substack{f\ne0\\\sum_u\pi_uf_u=0}}
 \frac{\sum_{uv\in E(G)}x_ux_v(f_u-f_v)^2}
 {\sum_u\pi_uf_u^2}.
\end{equation}
For weighted mean-zero $f$, direct expansion also gives
\begin{equation}\label{pre:variance}
 \sum_u\pi_uf_u^2=\sum_{u<v}\pi_u\pi_v(f_u-f_v)^2,
\end{equation}
where any fixed ordering of the vertices can be used.

\begin{lemma}\label{pre:diameter}
If $G$ is connected, has order $n\ge3$ and diameter $D$, then
\[
 \delta(G)\ge\frac{2}{D(n-1)\lambda_1^{D-1}}.
\]
\end{lemma}
\begin{proof}
For each unordered pair $u,v$, choose a shortest path $Q_{uv}$ and put
\[
 R_{uv}=\sum_{ab\in E(Q_{uv})}\frac1{x_ax_b}.
\]
Cauchy--Schwarz on this path gives
$(f_u-f_v)^2\le R_{uv}\mathcal E(f)$, where $\mathcal E(f)$ is the
numerator in \eqref{pre:variational}. If $ab$ is oriented toward $v$ on
an $h$-edge $u$--$v$ path, successive Perron equations imply
$x_ux_v/(x_ax_b)\le\lambda_1^{h-1}$. Since $\lambda_1\ge1$,
\[
 \sum_{u<v}x_u^2x_v^2R_{uv}
 \le D\lambda_1^{D-1}\sum_{u<v}x_ux_v
 \le\frac{D(n-1)\lambda_1^{D-1}}2.
\]
The last inequality follows from $\norm{\mathbf x}_1^2\le n$.
Now apply \eqref{pre:variance} and \eqref{pre:variational}.
\end{proof}

\subsection{Exact path identities}
Define $\phi_j(z)=\det(zI-A(P_j))$ for $j\ge0$, with $\phi_0=1$ and
$\phi_{-1}=0$. Expansion at a path endpoint gives
$\phi_j=z\phi_{j-1}-\phi_{j-2}$. For $z>2$, write
\[
 r(z)=\frac{z-\sqrt{z^2-4}}2\in(0,1),\qquad z=r(z)+r(z)^{-1}.
\]
\begin{lemma}\label{pre:path}
For $j\ge0$ and $z>2$,
\[
 \phi_j(z)=\frac{r(z)^{-j-1}-r(z)^{j+1}}{r(z)^{-1}-r(z)},
 \qquad 0<\phi_j(z)\le z^j.
\]
For a path with $L\ge2$ edges, the endpoint diagonal entries and the opposite-corner entry
of its internal resolvent $(zI-A(P_{L-1}))^{-1}$ are
\begin{align}
 s_L(z)&=\frac{\phi_{L-2}(z)}{\phi_{L-1}(z)}
 =r(z)\frac{1-r(z)^{2L-2}}{1-r(z)^{2L}},\label{pre:s}\\
 h_L(z)&=\frac1{\phi_{L-1}(z)}
 =\frac{(1-r(z)^2)r(z)^{L-1}}{1-r(z)^{2L}}.\label{pre:h}
\end{align}
In particular, $s_L(z)-r(z)=-h_L(z)r(z)^L$.
\end{lemma}
\begin{proof}
The formula for $\phi_j$ satisfies its recurrence and initial values.
Its positivity also follows from $z>\rho(P_j)$; the upper bound follows
inductively from the recurrence. Cofactor expansion gives the two stated
resolvent entries. Substitution of the formula for $\phi_j$ proves
\eqref{pre:s}--\eqref{pre:h} and their final identity.
\end{proof}
The expression in \eqref{pre:h} will also be denoted by $h_1(z)=1$ when
$L=1$; in that case it represents a direct edge.

\begin{lemma}[Centering and exact resistance]\label{pre:center}
There exists $c\in\mathbb R$ such that $\mathbf w=\mathbf y-c\mathbf x$
satisfies $\max_u w_u=M=-\min_u w_u>0$. It has the properties
\begin{equation}\label{pre:centered}
 \norm{\mathbf w}^2=1+c^2,\quad M^2\ge\frac1n,\quad
 |c|/M\le\sqrt n,\quad (\lambda_2I-A)\mathbf w=c\delta\mathbf x.
\end{equation}
If $w_u>0>w_v$ and $h=\dist(u,v)$, then
\begin{equation}\label{pre:exact-resistance}
 \delta(G)\ge\frac{4w_u|w_v|}{\phi_{h-1}(\lambda_1)}.
\end{equation}
\end{lemma}
\begin{proof}
The continuous function
$c\mapsto\max_u(y_u-cx_u)+\min_u(y_u-cx_u)$ is strictly decreasing
and tends to opposite infinities at the two ends of the real line. Its
zero gives the asserted centering. Here $M>0$, since otherwise
$\mathbf y$ would be a multiple of $\mathbf x$. Orthogonality gives
$\norm{\mathbf w}^2=1+c^2$ and $\mathbf x^T\mathbf w=-c$.
Thus $nM^2\ge1$ and $|c|\le M\norm{\mathbf x}_1\le M\sqrt n$.
The last equation in \eqref{pre:centered} is obtained by applying
$\lambda_2I-A$.

Choose a shortest path $u=v_0,\ldots,v_h=v$. Set $a_0=x_u$, $a_h=x_v$
and solve $\lambda_1a_i=a_{i-1}+a_{i+1}$ in its interior. The inverse
of $\lambda_1I-A(P_{h-1})$ is nonnegative: the internal path is a proper
principal subgraph of the connected graph, so its spectral radius is
strictly below $\lambda_1$. The Perron equations, with all contributions
off the path retained as nonnegative forcing, imply $x_{v_i}\ge a_i>0$.
Put $b_0=0$ and $b_i=\phi_{i-1}(\lambda_1)$ for $i\ge1$.
The shared recurrence implies
\[
 a_ib_{i+1}-a_{i+1}b_i=x_u\qquad(0\le i<h).
\]
Divide by $a_ia_{i+1}$ and telescope to obtain
\begin{equation}\label{pre:exact-path-resistance}
 \sum_{i=0}^{h-1}\frac1{x_{v_i}x_{v_{i+1}}}
 \le\sum_{i=0}^{h-1}\frac1{a_ia_{i+1}}
 =\frac{\phi_{h-1}(\lambda_1)}{x_ux_v}.
\end{equation}
This argument includes $h=1$, with no interior equations.
By \eqref{pre:energy} applied to $\mathbf y$,
\[
 \delta=\sum_{ab\in E(G)}x_ax_b
 \left(\frac{y_a}{x_a}-\frac{y_b}{x_b}\right)^2
 =\sum_{ab\in E(G)}x_ax_b
 \left(\frac{w_a}{x_a}-\frac{w_b}{x_b}\right)^2.
\]
Cauchy--Schwarz along the path and
\eqref{pre:exact-path-resistance} give
\[
 \delta\ge
 \frac{(w_ux_v+|w_v|x_u)^2}{x_ux_v\phi_{h-1}(\lambda_1)}
 \ge\frac{4w_u|w_v|}{\phi_{h-1}(\lambda_1)}.
\]
\end{proof}

\subsection{A two-dimensional Schur-complement estimate}
We record the elementary matrix argument used in both proofs.
\begin{lemma}\label{pre:schur}
Let $A$ be the adjacency matrix of a connected graph of order at least
three. Partition its
vertices into two selected vertices and the remaining vertices, so that
\[
 A=\begin{pmatrix}D&W^T\\ W&B\end{pmatrix},\qquad
 F(z)=zI_2-D-W^T(zI-B)^{-1}W.
\]
Suppose $\lambda_2>\rho(B)$ and
$I_2\preceq F'(z)\preceq aI_2$ for $\lambda_2\le z\le\lambda_1$.
Writing $F_{12}(z)=-g(z)$, one has
\[
 0\le F_{ii}(\lambda_1)\le a\delta,
 \qquad \delta\ge\frac{2g(\lambda_1)}a.
\]
\end{lemma}
\begin{proof}
Interlacing gives $\lambda_3\le\rho(B)<\lambda_2$, so $\lambda_2$
is simple. Block elimination is a congruence between $zI-A$ and
$\operatorname{diag}(F(z),zI-B)$. Thus the two eigenvalues of
$F(\lambda_2)$ are negative and zero, while those of $F(\lambda_1)$
are zero and positive. Integrating the upper derivative bound and using
the variational principle gives
$\lambda_{\max}(F(\lambda_1))\le a\delta$.
The matrix $F(\lambda_1)$ is positive semidefinite and singular.
Its diagonal entries are therefore nonnegative and at most $a\delta$;
their product is $g(\lambda_1)^2$, so its trace is at least
$2g(\lambda_1)$. Here $g(z)>0$: a direct edge or a walk through the
remaining vertices connects the two selected vertices, and the resolvent
has a nonnegative walk expansion. The trace bound proves the result.
Finally, differentiation gives
$F'(z)=I_2+W^T(zI-B)^{-2}W$, explaining the lower derivative bound.
\end{proof}

\section{Two comparison families}\label{sec:benchmarks}

\begin{lemma}[Double comets]\label{bench:comet}
If $k\ge3$, $L\ge1$, and $r=(k-1)^{-1/2}$, then
\begin{equation}\label{bench:comet-bound}
 \delta(\Comet(k,L))\le
 \frac{2(1-r^2)r^{L-1}}{1-r^{2L}}
 \le2(k-1)^{-(L-1)/2}.
\end{equation}
\end{lemma}
\begin{proof}
Put $\alpha=r+r^{-1}=k/\sqrt{k-1}$. Number the joining path
$0,\ldots,L$. Give a vector $\mathbf u$ the value $r^j$ at vertex $j$
and the value of its adjacent endpoint divided by $\alpha$ at each leaf.
Let $\mathbf v$ be its reflection. Substituting into the adjacency
equations gives
\[
 (A-\alpha I)\mathbf u=e\mathbf e_L,\quad
 (A-\alpha I)\mathbf v=e\mathbf e_0,\quad
 e=(1-r^2)r^{L-1}.
\]
The positive vector $\mathbf u+\mathbf v$ implies
$\lambda_1\le\alpha+e/(1+r^L)$. The Perron vector is invariant under
reflection, so $\mathbf u-\mathbf v$ is orthogonal to it. Since
\[
 \norm{\mathbf u-\mathbf v}^2\ge2(1-r^L)^2,\qquad
 (\mathbf u-\mathbf v)^T(A-\alpha I)(\mathbf u-\mathbf v)
 =-2e(1-r^L),
\]
the variational principle gives $\lambda_2\ge\alpha-e/(1-r^L)$.
Subtracting proves the first bound. The second uses
$1-r^{2L}\ge1-r^2$.
\end{proof}

For $k\ge3$ define
\begin{equation}\label{bench:clique-parameters}
 \beta_k=\frac{k-2+\sqrt{k^2-4}}2,\qquad
 \kappa_k=\beta_k+\beta_k^{-1}.
\end{equation}
\begin{lemma}[Double kites]\label{bench:kite}
For $k\ge3$ and $L\ge1$,
\begin{equation}\label{bench:kite-bound}
 \delta(\Kite(k,L))\le\frac{32}{k+3}\beta_k^{-(L-1)},
 \qquad \kappa_k-(k-1)>\frac1{k(k-1)}.
\end{equation}
\end{lemma}
\begin{proof}
In the unit Perron vector of $\Kite(k,L)$, let $a$ be the common
coordinate at the two clique roots. Its maximum degree is $k$, so
$\lambda_1\le k$. Each nonroot clique coordinate equals
$a/(\lambda_1-k+2)\ge a/2$; consequently $a^2\le2/(k+3)$.

For completeness, $\beta_k$ satisfies
$\beta_k^2-(k-2)\beta_k-(k-2)=0$. At $z=\kappa_k$ this gives
\[
 z-\frac{k-1}{z-k+2}=\beta_k^{-1}.
\]
Use the path values $r^j$, where $r=\beta_k^{-1}$, and use at each
nonroot clique vertex its root value divided by $\kappa_k-k+2$.
The residual of $A-\kappa_kI$ is positive only at the far root.
Adding its reflection gives a positive vector with a nonzero nonnegative
residual. Multiplication by the Perron vector proves
$\lambda_1>\kappa_k>2$. Hence $r_1=r(\lambda_1)<\beta_k^{-1}$.
The exact Perron coordinates along the joining path are
\[
 x_j=a\frac{r_1^j+r_1^{L-j}}{1+r_1^L}
 \qquad(0\le j\le L).
\]
Apply \eqref{pre:variational} to a function taking values $+1$ and $-1$
on the two reflected halves. If $L$ is odd, only the central edge
contributes and its energy is at most $16a^2r_1^{L-1}$; the weighted
norm is one. If $L$ is even, give the central vertex value zero.
The energy of the two central edges is at most $8a^2r_1^{L-1}$.
Both neighboring Perron coordinates exceed the central coordinate,
because $\lambda_1>2$; thus the squared weighted norm is at least $2/3$.
In either case the weighted mean is zero by reflection. These estimates,
$a^2\le2/(k+3)$, and $r_1<\beta_k^{-1}$ prove the gap bound.

Finally set $\varepsilon=k-1-\beta_k$. Direct substitution gives
$\varepsilon^2-k\varepsilon+1=0$, so $\varepsilon>1/k$.
The same identity gives
$\kappa_k-(k-1)=\varepsilon/\beta_k>1/[k(k-1)]$.
\end{proof}

\section{The spectral scale and concentration of branching in trees}
\label{tree:scale-section}

Consider an arbitrary sequence of trees attaining $\tau_n$, with their
orders tending to infinity. In this and the following two sections, every
asymptotic statement is along this sequence. We suppress the argument $T$
in its eigenvalues, write $\delta=\tau_n$, and set
\[
 D=\operatorname{diam}(T),\qquad m=n-D-1,\qquad
 a_n=\log n-\log\log n-\log 2.
\]
Thus $m$ is the number of vertices outside any diametral path. The
arbitrariness of the sequence will give a uniform eventual conclusion.

\begin{lemma}\label{tree:scales}
Along every such sequence,
\begin{align}
 -\log\delta&=\frac n2(a_n-1)+o(n),\label{tree:gap-scale}\\
 m&\sim\frac{n}{\log n},\label{tree:m-scale}\\
 \log\lambda_1&=\frac12a_n+o(1),\qquad
 \lambda_1^2\sim\lambda_2^2\sim\frac m2.\label{tree:eigen-scale}
\end{align}
In particular, $n^C\delta\longrightarrow0$ for every fixed $C>0$.
\end{lemma}

\begin{proof}
Choose $k=\lfloor n/(2\log n)\rfloor$ and $L=n-2k-1$.
The tree $\Comet(k,L)$ has order $n$, and Lemma~\ref{bench:comet}
therefore gives
\begin{equation}\label{tree:benchmark}
 -\log\delta\geq
 \left(\frac n2-k-1\right)\log(k-1)-\log2
 =\frac n2(a_n-1)+o(n).
\end{equation}
Indeed, $\log(k-1)=a_n+O(\log n/n)$ and
$k\log(k-1)=n/2+o(n)$. This already proves the last assertion.

Fix a diametral path $P$. Regard its adjacency matrix as indexed by
$V(T)$, with zero rows at vertices outside the path, and write
$A(T)=A(P)+A(F)$, where $F$ consists of the edges outside $P$.
The forest $F$ has exactly $m$ edges. Spectral symmetry of a bipartite
graph and $\operatorname{tr}A(F)^2=2m$ show that
\[
 \sum_{\lambda_j(F)>0}\lambda_j(F)^2=m,
 \qquad \lambda_2(F)\leq\sqrt{m/2}.
\]
For the second assertion, there is nothing to prove if
$\lambda_2(F)\leq0$; otherwise the two largest positive squared
eigenvalues are each at least $\lambda_2(F)^2$. Weyl's inequality now
gives
\[
 \lambda_2\leq\lambda_2(F)+\lambda_1(P)<\sqrt{m/2}+2.
\]
As $\delta<1$ eventually, we also have
$\lambda_1<\sqrt{m/2}+3$. Lemma~\ref{pre:diameter} consequently implies
\begin{equation}\label{tree:rough-upper}
 -\log\delta\leq
 (n-m-2)\log\bigl(3+\sqrt{m/2}\bigr)+O(\log n).
\end{equation}

First, $m\to\infty$: if $m$ were bounded on a subsequence, the
right-hand side would be $O(n)$, contradicting
\eqref{tree:benchmark}. Next, $m/n\to0$. Otherwise, on a subsequence
with $m\geq\varepsilon n$ for some $\varepsilon>0$, the coefficient of
$n\log n$ in \eqref{tree:rough-upper} would be at most
$(1-\varepsilon)/2$, again contradicting the benchmark. Since
\[
 \log\bigl(3+\sqrt{m/2}\bigr)
 =\frac12\log(m/2)+O(m^{-1/2}),
\]
equation~\eqref{tree:rough-upper} becomes
\[
 -\log\delta\leq\frac{n-m}{2}\log(m/2)+o(n).
\]
Here the discarded error is $O(n/\sqrt m+\log n)=o(n)$.

Put $c=m\log n/n$. Comparison with \eqref{tree:benchmark} yields
\begin{equation}\label{tree:c-stability}
 0\leq\log c-c+1+\frac mn\log\frac{2n}{m}+o(1).
\end{equation}
To see this directly, use $\log(m/2)=a_n+\log c$ and
$(m/n)\log(m/2)=c-(m/n)\log(2n/m)$. The penultimate term in
\eqref{tree:c-stability} tends to zero. The function
$\log c-c+1$ is nonpositive, vanishes only at $c=1$, and tends to
$-\infty$ as $c\to0$ or $c\to\infty$. Thus $c\to1$, proving
\eqref{tree:m-scale}.

Apply Lemma~\ref{pre:diameter} once more, now retaining $\lambda_1$:
\[
 \log\lambda_1\geq
 \frac{\frac n2(a_n-1)+o(n)}{n-m-2}
 =\frac12a_n+o(1).
\]
The last equality follows from
$(a_n-1)m/n\to1$ and $a_n(m/n)^2\to0$.
Conversely, $\lambda_1<\sqrt{m/2}+3$ and
\eqref{tree:m-scale} give $\log\lambda_1\leq a_n/2+o(1)$.
It follows that $\lambda_1^2\sim m/2$. The same holds for
$\lambda_2^2$, since $\lambda_1-\lambda_2=\delta$ is smaller than
every fixed inverse power of $n$. Finally,
\[
 -\log\delta\leq(n-m-2)\log\lambda_1+O(\log n)
 =\frac n2(a_n-1)+o(n),
\]
where we used $ma_n=n(1+o(1))$. This completes the proof.
\end{proof}

\begin{lemma}\label{tree:mass}
Let $H$ be a nontrivial tree with $\lambda_1(H)>2$, and let
$\mathbf z>0$ be a Perron eigenvector of $H$. If $H$ has $a$ leaves,
then
\[
 \bigl(\lambda_1(H)-2\bigr)\sum_{v\in V(H)}z_v
 \leq(a-2)\max_{v\in V(H)}z_v.
\]
\end{lemma}

\begin{proof}
Sum the eigenvalue equations and subtract twice the sum of the
coordinates. The resulting identity is
\[
 \bigl(\lambda_1(H)-2\bigr)\sum_vz_v
 =\sum_v(d_H(v)-2)z_v.
\]
Discard the negative terms, and use
$\sum_v(d_H(v)-2)_+=a-2$.
\end{proof}

\begin{lemma}\label{tree:hubs}
There are distinct vertices $u_1,u_2$ of degrees $d_1,d_2$ such that
\begin{equation}\label{tree:hub-scales}
 d_1\sim d_2\sim m/2.
\end{equation}
Moreover, if $B=A(T-\{u_1,u_2\})$, then
\begin{equation}\label{tree:residual-small}
 \lVert B\rVert=o(\sqrt m)=o(\lambda_2).
\end{equation}
\end{lemma}

\begin{proof}
The unit eigenvector $\mathbf y$ corresponding to $\lambda_2$ has
both signs, because it is orthogonal to the positive unit eigenvector
$\mathbf x$ corresponding to $\lambda_1$. Choose a connected
component $H_1$ of the subgraph induced by $\{v:y_v>0\}$ and a
connected component $H_2$ of the subgraph induced by $\{v:y_v<0\}$.
At a vertex of $H_1$, the omitted neighbors have nonpositive
coordinates, so
\[
 A(H_1)\mathbf y_{H_1}\geq\lambda_2\mathbf y_{H_1}.
\]
The analogous inequality holds for $-\mathbf y_{H_2}>0$.
Multiplying by positive Perron eigenvectors of the respective
components, and applying interlacing for the upper bound, gives
\begin{equation}\label{tree:nodal-radii}
 \lambda_2\leq\lambda_1(H_i)\leq\lambda_1
 \qquad(i=1,2).
\end{equation}
In particular, the components are nontrivial and have spectral radius
greater than two eventually.

The endpoints of a diametral path are leaves, and no internal vertex
of that path is a leaf. Thus $T$ has at most $m+2$ leaves. Every
nontrivial connected subtree $H$ has at most this many leaves, since
\[
 \#\{\text{leaves of }H\}-2
 =\sum_{v\in V(H)}(d_H(v)-2)_+
 \leq\sum_{v\in V(T)}(d_T(v)-2)_+\leq m.
\]
Let $\mathbf z>0$ be a Perron eigenvector of $H_i$, and choose
$u_i$ at a maximum coordinate $M$. By Lemma~\ref{tree:mass},
\[
 \frac{\sum_v z_v}{M}
 \leq\frac{m}{\lambda_1(H_i)-2}=O(\sqrt m).
\]
There is exactly one walk of length two from $u_i$ to each vertex
at distance two in a tree, and there are $d_{H_i}(u_i)$ such walks
back to $u_i$. Consequently,
\[
 \lambda_1(H_i)^2M
 =d_{H_i}(u_i)M+
   \sum_{\operatorname{dist}_{H_i}(u_i,v)=2}z_v.
\]
Together with \eqref{tree:nodal-radii}, this implies
\[
 d_i=d_T(u_i)\geq\lambda_2^2-O(\sqrt m)=m/2-o(m).
\]
The vertices $u_1,u_2$ are distinct. Their contributions to the
positive degree excess are $d_1-2$ and $d_2-2$ eventually, while the
total excess is at most $m$. Hence both degrees are asymptotic to
$m/2$, and the total positive degree excess at all remaining vertices
is $o(m)$. Their degrees are therefore uniformly $o(m)$.

For completeness, a forest of maximum degree $\Delta\geq1$ has
operator norm at most $2\sqrt\Delta$. Orient each component away
from a root. For an edge with parent $a$ and child $b$, use
\[
 2|z_a z_b|\leq\Delta^{-1/2}z_a^2+\Delta^{1/2}z_b^2.
\]
Each vertex is a parent at most $\Delta$ times and a child at most
once, so summation bounds the absolute adjacency quadratic form by
$2\sqrt\Delta\lVert\mathbf z\rVert^2$. An edgeless forest has norm
zero. Apply this estimate to $T-\{u_1,u_2\}$, whose maximum degree
is $o(m)$, to obtain \eqref{tree:residual-small}.
\end{proof}

\section{Equal degrees at the two dominant vertices}
\label{tree:degree-section}

Retain $u_1,u_2$ and $B$ from Lemma~\ref{tree:hubs}. Let
$\mathbf a_i$ be the indicator vector of the neighbors of $u_i$
in $T-\{u_1,u_2\}$, let $W=[\mathbf a_1\ \mathbf a_2]$, and let
$A_U=A(T[\{u_1,u_2\}])$. For $z>\lVert B\rVert$, define
\begin{equation}\label{tree:schur-definition}
 S(z)=W^{\mathsf T}(zI-B)^{-1}W,
 \qquad F(z)=zI_2-A_U-S(z).
\end{equation}
We write $F_{12}(z)=-g(z)$; explicitly,
\[
 g(z)=\mathbf1_{\{u_1u_2\in E(T)\}}
      +\mathbf a_1^{\mathsf T}(zI-B)^{-1}\mathbf a_2>0.
\]
The strict inequality follows from the unique path between the two
vertices and the nonnegative Neumann expansion of the resolvent.

\begin{lemma}\label{tree:schur-bounds}
Eventually, throughout $\lambda_2\leq z\leq\lambda_1$,
\[
 I_2\preceq F'(z)\preceq4I_2.
\]
Consequently,
\begin{equation}\label{tree:schur-consequences}
 \delta\geq\frac12g(\lambda_1),\qquad
 0\leq F_{ii}(\lambda_1)\leq4\delta\quad(i=1,2).
\end{equation}
\end{lemma}

\begin{proof}
By \eqref{tree:residual-small}, $\lambda_2>\lVert B\rVert$
eventually. Interlacing gives $\lambda_3\leq\lambda_1(B)<\lambda_2$,
where $\lambda_1(B)$ denotes the largest eigenvalue of the matrix
$B$. Thus the two-root Schur complement in Lemma~\ref{pre:schur}
applies. Direct differentiation gives
\[
 F'(z)=I_2+W^{\mathsf T}(zI-B)^{-2}W.
\]
The second term is positive semidefinite, and its norm is at most
\[
 \frac{\lVert W\rVert^2}{(z-\lVert B\rVert)^2}
 \leq\frac{d_1+d_2}{(\lambda_2-\lVert B\rVert)^2}
 =2+o(1).
\]
The upper bound $4I_2$ follows eventually. Apply
Lemma~\ref{pre:schur} with derivative bound four to obtain
\eqref{tree:schur-consequences}.
\end{proof}

\begin{lemma}\label{tree:budget}
Let $L=\operatorname{dist}_T(u_1,u_2)$. Then
\begin{equation}\label{tree:separation}
 L\geq n-(1+o(1))\frac{n}{\log n}.
\end{equation}
In particular, $L\geq3$ eventually. The integer
\begin{equation}\label{tree:b-definition}
 b=n-(L+1)-(d_1-1)-(d_2-1)
\end{equation}
is nonnegative and satisfies $b=o(m)$.
\end{lemma}

\begin{proof}
For $L\geq2$, the unique connecting path contributes the term
$z^{-(L-1)}$ to the Neumann expansion of $g(z)$. For $L=1$ the
direct edge gives $g(z)\geq1$, so in all cases
$g(\lambda_1)\geq\lambda_1^{-(L-1)}$. By
\eqref{tree:schur-consequences},
\[
 (L-1)\log\lambda_1\geq-\log\delta-\log2.
\]
Equations~\eqref{tree:benchmark} and \eqref{tree:eigen-scale} give
\[
 L-1\geq n\frac{a_n-1+o(1)}{a_n+o(1)}
 =n-(1+o(1))\frac{n}{\log n}.
\]
The additive $o(1)$ precision in $\log\lambda_1$ is used here.

The path has $L+1$ vertices. At its endpoints there are $d_1-1$
and $d_2-1$ neighbors outside the path. These neighbors are all
distinct, by acyclicity, so $b\geq0$. Finally,
\eqref{tree:m-scale}, \eqref{tree:hub-scales}, and
\eqref{tree:separation} yield
\[
 0\leq b=n-L+1-d_1-d_2=o(n/\log n)=o(m).
\]
\end{proof}

\begin{lemma}\label{tree:diagonal}
For $L\geq3$ and $z>\lVert B\rVert$,
\begin{equation}\label{tree:diagonal-estimate}
 S_{ii}(z)=\frac{d_i}{z}+E_i(z),\qquad
 0\leq E_i(z)\leq\frac{b+1}{z(z^2-\lVert B\rVert^2)}.
\end{equation}
\end{lemma}

\begin{proof}
Since the hubs are not adjacent, $\lVert\mathbf a_i\rVert^2=d_i$.
Each entry of $B\mathbf a_i$ is zero or one: two neighbors in the
support of $\mathbf a_i$ would create a cycle of length four through
$u_i$. Its nonzero entries therefore count vertices at distance two
from $u_i$ reached without passing through the other hub.
Exactly one such vertex lies on the connecting path. Any other such
vertex lies outside that path and cannot be an immediate off-path
neighbor of either hub. For its own hub this would create a triangle;
for the other hub it would create a second path between the hubs and
hence a cycle. Thus every such vertex is among the $b$ vertices
counted by \eqref{tree:b-definition}, and
\begin{equation}\label{tree:second-neighborhood}
 \lVert B\mathbf a_i\rVert^2\leq b+1.
\end{equation}

All neighbors of $u_i$ lie in the same bipartition class of $T$.
Consequently $\mathbf a_i^{\mathsf T}B^{2j+1}\mathbf a_i=0$ for
every $j\geq0$. Expanding the resolvent gives
\[
 S_{ii}(z)=\frac{d_i}{z}
 +\sum_{j\geq1}\frac{\lVert B^j\mathbf a_i\rVert^2}{z^{2j+1}}.
\]
For $j\geq1$, the numerator is at most
$\lVert B\rVert^{2j-2}\lVert B\mathbf a_i\rVert^2$.
Summing the resulting geometric series proves the claim.
\end{proof}

\begin{lemma}\label{tree:equal-degrees}
Eventually $d_1=d_2=k+1$ for an integer $k$, and
\begin{equation}\label{tree:k-scale}
 k\sim\frac{n}{2\log n},\qquad b=o(k),\qquad
 n=2k+L+1+b.
\end{equation}
Moreover,
\begin{equation}\label{tree:radius-upper}
 \lambda_1^2\leq k+1+\frac{3(b+1)}{k}.
\end{equation}
\end{lemma}

\begin{proof}
From the definitions, there is the exact identity
\[
 \lambda_1^2-d_i
 =\lambda_1E_i(\lambda_1)+\lambda_1F_{ii}(\lambda_1).
\]
Use \eqref{tree:diagonal-estimate} and
\eqref{tree:schur-consequences} to obtain
\begin{equation}\label{tree:integer-interval}
 d_i\leq\lambda_1^2\leq
 d_i+\frac{b+1}{\lambda_1^2-\lVert B\rVert^2}
       +4\lambda_1\delta.
\end{equation}
The error on the right tends to zero: $b=o(m)$,
$\lambda_1^2-\lVert B\rVert^2\sim m/2$, and
$\lambda_1\delta\to0$. The integers $d_1,d_2$ therefore lie in the
same interval of length $o(1)$ and must be equal eventually.
Their common value and the definition of $b$ give
\eqref{tree:k-scale}.

For large $n$,
$\lambda_1^2-\lVert B\rVert^2\geq k/2$.
Also $4\lambda_1\delta\leq(b+1)/k$, because $b+1\geq1$,
$k\lambda_1$ grows at most polynomially in $n$, and $\delta$ is
smaller than every fixed inverse power of $n$. Substitution in
\eqref{tree:integer-interval} proves \eqref{tree:radius-upper}.
\end{proof}

\section{Exact completion of the tree structure}
\label{tree:completion-section}

\begin{lemma}\label{tree:path-comparison}
Write $\lambda_1=r+r^{-1}$ with $0<r<1$, and put $t=r^{-2}$.
Along the minimizing sequence, eventually
\begin{equation}\label{tree:path-comparison-bound}
 \delta\geq\frac14t^{-(L-1)/2},\qquad
 t\leq k-1+\frac{3(b+1)}k.
\end{equation}
\end{lemma}

\begin{proof}
The representation of $\lambda_1$ is valid eventually because
$\lambda_1\to\infty$. Retain in the residual resolvent only those
walks confined to the interior of the path between $u_1$ and $u_2$.
All terms are nonnegative, so Lemma~\ref{pre:path} gives
\[
 g(\lambda_1)\geq
 \bigl((\lambda_1I-A(P_{L-1}))^{-1}\bigr)_{1,L-1}
 =\frac1{\phi_{L-1}(\lambda_1)}
 =\frac{(1-r^2)r^{L-1}}{1-r^{2L}}.
\]
Since $r\to0$, the last expression is at least
$r^{L-1}/2$ eventually. The first assertion now follows from
$\delta\geq g(\lambda_1)/2$.
For the second, use
$\lambda_1^2=t+2+t^{-1}$ and \eqref{tree:radius-upper}, discarding
the positive term $t^{-1}$.
\end{proof}

\begin{proof}[Proof of Theorem~\ref{thm:trees}]
Suppose that there are minimizers which are not symmetric double
comets at arbitrarily large orders. Choose a sequence of these
minimizers and apply the preceding lemmas. The symmetric double
comet
\[
 C=\Comet(k,L+b)
\]
has the same order $n=2k+L+1+b$. By Lemmas~\ref{bench:comet} and
\ref{tree:path-comparison},
\begin{equation}\label{tree:gap-ratio}
 \frac{\delta(T)}{\delta(C)}
 \geq\frac18(k-1)^{b/2}
 \left(1+\frac{3(b+1)}{k(k-1)}\right)^{-(L-1)/2}.
\end{equation}
Taking logarithms and using $\log(1+s)\leq s$ for $s\geq0$ yields
\begin{equation}\label{tree:log-ratio}
 \log\frac{\delta(T)}{\delta(C)}
 \geq\frac b2\log(k-1)
       -\frac{3(L-1)(b+1)}{2k(k-1)}-\log8.
\end{equation}
By \eqref{tree:k-scale}, $L\leq n=O(k\log k)$.
If $b\geq1$, then $b+1\leq2b$, and the right-hand side is bounded
below by
\[
 b\left(\frac12\log(k-1)-O\left(\frac{\log k}{k}\right)\right)
 -\log8.
\]
This tends to $+\infty$ uniformly over integers $b\geq1$,
contradicting the minimality of $T$.

It follows that $b=0$ eventually. Every vertex then lies on the
connecting path or is an immediate off-path neighbor of one of its
endpoints. Acyclicity forces each such off-path neighbor to be a
leaf: the path together with its incident edges to these neighbors
already forms a spanning tree. The equality $d_1=d_2=k+1$ gives
exactly $k$ leaves at each endpoint. Hence $T=\Comet(k,L)$,
contradicting our choice of the sequence. This proves the existence
of a uniform eventual threshold. Finally,
$k\sim n/(2\log n)$ and $L=n-2k-1$ give the stated parameter
asymptotics.
\end{proof}

\section{The global scale and vertex budget for graphs}
\label{sec:graph-scale}

We now consider an arbitrary sequence $G\in\mathcal G_n$ with
$\delta(G)=\eta_n$ and $n\to\infty$. Use the centered vector
$\mathbf w=\mathbf y-c\mathbf x$ of Lemma~\ref{pre:center}, and choose
vertices $a,b$ with $w_a=M=-w_b$. Write $d=\dist(a,b)$ and
\begin{equation}\label{graph:k0}
 k_0=\left\lceil\lambda_2-\sqrt n\,\delta\right\rceil.
\end{equation}
The following counting bound is valid before imposing minimality,
whenever $k_0\ge1$.

\begin{lemma}\label{graph:neighborhoods}
One has $d\le n-2k_0+1$. If also $\lambda_1>2$, then
\begin{equation}\label{graph:universal}
 \delta\ge\frac4n\lambda_1^{-n+2\lambda_1-2(1+\sqrt n)\delta}.
\end{equation}
\end{lemma}
\begin{proof}
At the positive extremum, \eqref{pre:centered} gives
\[
 \lambda_2M=\sum_{v\sim a}w_v+c\delta x_a
 \le M|N(a)\cap\{w>0\}|+|c|\delta x_a.
\]
Since $x_a\le1$ and $|c|/M\le\sqrt n$, at least $k_0$ neighbors of
$a$ have positive centered coordinates. The negative extremum has at
least $k_0$ negative neighbors. These two sets are disjoint and each
meets a shortest $a$--$b$ path in at most one vertex. Their union with
that path has at least $d+1+2k_0-2$ vertices. This proves the distance
bound and, by \eqref{graph:k0},
\[
 d-1\le n-2\lambda_1+2(1+\sqrt n)\delta.
\]
Lemma~\ref{pre:center} at $a,b$ gives
$\delta\ge4M^2/\phi_{d-1}(\lambda_1)$. Now use $M^2\ge1/n$ and
$\phi_{d-1}(\lambda_1)\le\lambda_1^{d-1}$.
\end{proof}

\begin{lemma}\label{graph:scale}
Let $z_n>1$ be the solution of $z_n(1+\log z_n)=n/2$, defined for
all sufficiently large $n$. Then
\begin{align}
 -\log\eta_n&=(n-2z_n)\log z_n+O(\log n),\label{graph:rate}\\
 \lambda_1&=z_n+O\!\left(\sqrt{n/\log n}\right),
 \qquad k_0\sim\frac{n}{2\log n}.\label{graph:scale-eq}
\end{align}
In particular, $\delta$ is smaller than every fixed inverse power of $n$
eventually.
\end{lemma}
\begin{proof}
The function $f_n(z)=(n-2z)\log z$, $z\ge1$, has a unique maximum at
$z_n$. Take an integer $h$ nearest to $z_n+2$, and compare with
$\Kite(h,n-2h+1)$. Since $\beta_h>h-2$, Lemma~\ref{bench:kite} gives
\[
 -\log\eta_n\ge(n-2h)\log(h-2)-O(\log n)
 \ge f_n(z_n)-O(\log n).
\]
To justify the second estimate, put $u=h-2=z_n+O(1)$.
The expression on the left is $f_n(u)-4\log u$; Taylor's theorem at
$z_n$ gives $f_n(u)=f_n(z_n)+O((\log n)^2/n)$.
This lower bound proves the asserted rapid decay of $\delta$.

Also $\lambda_1\to\infty$. Otherwise, along a subsequence on which it
is bounded, Lemma~\ref{pre:diameter}, with $D\le n-1$, would give
$-\log\delta=O(n+\log n)$, a contradiction. Thus $\lambda_1>2$ and
$k_0\ge1$ eventually. Applying \eqref{graph:universal}, and using
$\lambda_1\le n-1$ and the rapid decay of $\delta$, gives
\[
 -\log\eta_n\le f_n(\lambda_1)+O(\log n)
 \le f_n(z_n)+O(\log n).
\]
This proves \eqref{graph:rate}. If $t=\lambda_1/z_n$, direct algebra gives
\[
 f_n(z_n)-f_n(tz_n)
 =2z_n\bigl[\log z_n(t-1-\log t)+(t-1)\log t\bigr].
\]
The two bracketed terms are nonnegative. Since their sum multiplied by
$2z_n$ is $O(\log n)$, first $t\to1$, and then
$(t-1)^2=O(\log n/n)$. Because $z_n\sim n/(2\log n)$, the claimed
bound on $\lambda_1$ follows. The definition of $k_0$ changes
$\lambda_1$ by only $O(1)$, proving its scale.
\end{proof}

\begin{lemma}[A constant vertex budget]\label{graph:slack}
Eventually the nonnegative integer
\[
 b_0=n-d-2k_0+1
\]
belongs to $\{0,1,2\}$.
\end{lemma}
\begin{proof}
Nonnegativity follows from Lemma~\ref{graph:neighborhoods}. Put
$h=k_0+1$, $B_0=\beta_h$, and $m=d-1=n-2k_0-b_0$.
The definition of $k_0$ and \eqref{bench:kite-bound} give
\[
 \lambda_1\le k_0+(1+\sqrt n)\delta<\kappa_h
\]
eventually; the separation of $\kappa_h$ from $k_0$ is at least
$1/[k_0(k_0+1)]$, whereas the error is superpolynomially small.
For $\gamma=(\lambda_1+\sqrt{\lambda_1^2-4})/2$, eventually
$2\le\gamma<B_0$. Lemma~\ref{pre:path} implies
\[
 \phi_m(\lambda_1)
 =\frac{\gamma^{m+1}-\gamma^{-m-1}}{\gamma-\gamma^{-1}}
 \le\frac43 B_0^m.
\]
The exact resistance bound at the extrema yields
$\delta\ge3B_0^{-m}/n$. On the other hand, the same-order competitor
$\Kite(k_0+1,n-2k_0-1)$ gives
\[
 \delta\le\frac{32}{k_0+4}B_0^{-(n-2k_0-2)}.
\]
Consequently $B_0^{b_0-2}\le32n/[3(k_0+4)]$. If $b_0\ge3$, this
forces $B_0(k_0+4)\le32n/3$, contrary to
$B_0\sim k_0\sim n/(2\log n)$.
\end{proof}

\section{Dense terminal cores in a minimizing graph}
\label{sec:graph-cores}

Fix $\theta=0.99$ and a shortest path $Q=(a=v_0,\ldots,v_d=b)$.
Define
\[
 A_+=N(a)\cap\{w>0\},\qquad B_-=N(b)\cap\{w<0\},\qquad
 E=V(G)\setminus(V(Q)\cup A_+\cup B_-).
\]
Let $I_A$ be the indicator of $v_1\in A_+$ and $I_B$ that of
$v_{d-1}\in B_-$. Inclusion--exclusion gives the exact identity
\begin{equation}\label{graph:count}
 b_0=(|A_+|-k_0)+(|B_-|-k_0)+(2-I_A-I_B)+|E|.
\end{equation}
Each term is nonnegative. Thus $|A_+|,|B_-|\le k_0+2$ and $|E|\le2$.
We use several consequences of shortestness. There is no edge from
$A_+$ to $B_-$, since such an edge would give $d\le3$.
A vertex of $A_+$ can have path neighbors only at positions $0,1,2$;
the reversed assertion holds for $B_-$. In general, the path neighbors
of any vertex outside a geodesic have indices spanning at most two:
otherwise the two incident edges would shorten the geodesic. In
particular there are at most three such neighbors.

\begin{lemma}\label{graph:cores}
Let
\[
 C_1=\{v:w_v\ge\theta M\},\qquad
 C_2=\{v:w_v\le-\theta M\}.
\]
Then, along the minimizing sequence:
\begin{enumerate}
\item $|C_i|=k_0+O(1)$, and the minimum degree of $G[C_i]$ is at least
$0.9|C_i|-1$ eventually;
\item every vertex of $C_i$ is within distance two of its corresponding
extremum, and
\begin{equation}\label{graph:core-deficits}
 \sum_{v\in C_1}(1-w_v/M)=O(1),\qquad
 \sum_{v\in C_2}(1+w_v/M)=O(1);
\end{equation}
\item the residual graph $G-(C_1\cup C_2)$ has bounded maximum degree
and bounded adjacency norm. Only $O(1)$ residual vertices meet the cores,
and every residual coordinate has absolute value less than $\theta M$.
\end{enumerate}
\end{lemma}
\begin{proof}
Discard the negative neighbor contributions in the equation at $a$.
This gives
\begin{equation}\label{graph:neighbor-deficit}
 \sum_{v\in A_+}(1-w_v/M)
 \le |A_+|-\lambda_2+\sqrt n\,\delta\le3+o(1).
\end{equation}
Since $1-\theta$ is fixed, only $O(1)$ vertices of $A_+$ fail the
threshold defining $C_1$. The analogous assertion holds at $b$.

For $3\le j\le d-3$, the vertex $v_j$ has at most its two path
neighbors and the vertices of $E$ as neighbors. Hence its degree is at
most four, and \eqref{pre:centered} yields
\[
 |w_{v_j}|\le\frac{4+\sqrt n\,\delta}{\lambda_2}M<\theta M.
\]
A high positive path vertex cannot be among the last three vertices of
$Q$: its positive neighbors would then be confined to a bounded set.
Similarly a high negative vertex cannot be at the positive end.
A high positive vertex in $E$ needs at least
$\theta\lambda_2-o(1)$ positive neighbors. Outside $A_+$ it has at
most three path neighbors and one other neighbor in $E$, and no positive
neighbor in $B_-$. Thus it meets $A_+$ and lies within distance two of
$a$. Reverse signs for the assertion about $C_2$.

It follows that $C_1$ consists of the high vertices of $A_+$, the
extremum $a$, and at most a bounded number of further path or $E$
vertices. This proves its order estimate, its localization, and its
deficit bound from \eqref{graph:neighbor-deficit}. The same holds for
$C_2$. Every vertex of $C_1$ needs
$\theta\lambda_2-o(1)=\theta k_0-O(1)$ positive neighbors. All but
$O(1)$ of these lie in $C_1$: the remaining possibilities are the
omitted vertices of $A_+$, a bounded number of path vertices, and $E$.
The opposite terminal set is too far away to contribute. Hence its
internal degree is at least $\theta k_0-O(1)$, which is at least
$0.9|C_1|-1$ eventually. The other core is identical in this respect.

After both cores are removed, only $O(1)$ vertices outside $Q$ remain.
Each has at most three neighbors on $Q$, so the residual graph has
bounded maximum degree and therefore bounded norm. The path neighbors
of a core vertex are confined to the first or last four positions,
because that core vertex is within distance two of the corresponding
extremum. Thus only boundedly many residual vertices meet either core.
Finally, the cores were defined using \emph{all} coordinates at the
threshold, so every residual coordinate is strictly below it in absolute
value.
\end{proof}

\section{Transmission and attachment estimates}
\label{sec:graph-attachments}

Interchange the two ends if necessary so that $|C_1|=k$ and
$|C_2|=k+q$, where $q\ge0$. When doing so replace
$(\mathbf w,\mathbf y,c)$ by $(-\mathbf w,-\mathbf y,-c)$, leaving
$\mathbf x$ unchanged. Thus $\mathbf w|_{C_1}>0$ and
$\mathbf w|_{C_2}<0$. By Lemma~\ref{graph:cores}, $q=O(1)$,
$k=k_0+O(1)$, and $\lambda_1,\lambda_2=k+O(1)$.

Let $\mathcal R=V(G)\setminus(C_1\cup C_2)$,
$A_i=A(G[C_i])$, and $B=A(G[\mathcal R])$.
Let $W_i$ be the matrix with rows indexed by $C_i$ and columns by
$\mathcal R$ that records the edges between these sets. Only boundedly
many columns are nonzero; consequently $\norm{W_i}^2=O(k)$.
Write $\mathbf w_i,\mathbf x_i$ for restrictions to $C_i$, and use
subscript $\mathcal R$ for restrictions to the residual set. Put
$R_2=(\lambda_2I-B)^{-1}$, which is entrywise nonnegative and has norm
$O(1/k)$.

Choose a shortest path between $C_1$ and $C_2$, of $L$ edges. Its first
and last internal vertices are $s_1,s_2$; let $d_i\ge1$ count the
neighbors of $s_i$ in its adjacent core. Define
\begin{equation}\label{graph:exception-count}
 b=n-|C_1|-|C_2|-(L-1)\ge0.
\end{equation}
Core localization gives $d-4\le L\le d$. Therefore
\begin{equation}\label{graph:corridor-scale}
 b=O(1),\qquad L=n-2k+O(1),\qquad
 \frac{L}{k\log k}\longrightarrow2.
\end{equation}
No internal path vertex other than $s_1,s_2$ meets a core, since that
would shorten the path. No residual vertex meets both cores. Both
assertions hold eventually because $L\to\infty$.

\begin{lemma}[Endpoint transmission]\label{graph:transmission}
Along every minimizing sequence,
\begin{equation}\label{graph:transmission-bound}
 \delta\ge\frac2k(\theta^2-o(1))d_1d_2\lambda_1^{-(L-1)}.
\end{equation}
\end{lemma}
\begin{proof}
The residual part of \eqref{pre:centered} is exactly
\begin{equation}\label{graph:residual-equation}
 \mathbf w_{\mathcal R}
 =R_2(W_1^T\mathbf w_1+W_2^T\mathbf w_2+c\delta\mathbf x_{\mathcal R}).
\end{equation}
The source from the cores has bounded support and coordinates of
magnitude $O(kM)$, so its norm is $O(kM)$. The forcing is negligible
because $|c|/M\le\sqrt n$ and $\delta$ is superpolynomially small.
It follows that $\norm{\mathbf w_{\mathcal R}}=O(M)$. Thus
\[
 1+c^2=\norm{\mathbf w}^2\le(2k+O(1))M^2,
 \qquad M^2\ge\frac{1-o(1)}{2k}.
\]
The zeroth term of the Neumann series for $R_2$ contributes at least
$\theta d_1M/\lambda_2$ to $w_{s_1}$ from the positive core. The
opposite-core source has residual distance at least $L-2$ from $s_1$.
If there were a shorter residual path, adjoining one edge into each
core would contradict the choice of $L$. Its total contribution has
absolute value at most
\[
 O(kM)\frac{\norm B^{L-2}}
 {\lambda_2^{L-1}(1-\norm B/\lambda_2)}.
\]
This is smaller than $M$ times every fixed inverse power of $k$, since
$\norm B=O(1)$ and $L\sim2k\log k$. The forcing term in
\eqref{graph:residual-equation} has the same property. Therefore
\[
 w_{s_1}\ge(\theta-o(1))\frac{d_1M}{\lambda_2},\qquad
 -w_{s_2}\ge(\theta-o(1))\frac{d_2M}{\lambda_2}.
\]
In the full graph the distance between $s_1,s_2$ is exactly $L-2$;
a shorter path would again yield a shorter walk between the cores.
Apply \eqref{pre:exact-resistance}, use
$\phi_{L-3}(\lambda_1)\le\lambda_1^{L-3}$, and discard the factor
$\lambda_1^2/\lambda_2^2\ge1$. The resulting bound is
\eqref{graph:transmission-bound}.
\end{proof}

For $v\in\mathcal R$ let $d_i(v)=|N(v)\cap C_i|$, and define
\begin{equation}\label{graph:attachment-mass}
 \sigma_i=\sum_{v\in\mathcal R}\left(\frac{d_i(v)}k\right)^2.
\end{equation}
Only boundedly many terms are nonzero.

\begin{lemma}[A finite-rank spectral bound]\label{graph:pole}
Uniformly over residual vertices,
\begin{equation}\label{graph:density}
 d_i(v)/k\le\theta+O(1/k).
\end{equation}
Moreover,
\begin{equation}\label{graph:pole-bound}
 \lambda_1\le k-1+\min(\sigma_1,q+\sigma_2)+O(1/k).
\end{equation}
\end{lemma}
\begin{proof}
Each residual vertex has bounded residual degree and meets at most one
core. By the \emph{total} deficit estimate \eqref{graph:core-deficits},
the sum of absolute centered coordinates of its neighbors in that core
is at least $d_i(v)M-O(M)$. Its residual neighbors contribute only
$O(M)$ in absolute value. Hence the approximate eigenvector equation
implies
\[
 d_i(v)M-O(M)\le\lambda_2|w_v|+|c|\delta x_v
 \le\theta\lambda_2M+o(M).
\]
Since $\lambda_2=k+O(1)$, this proves \eqref{graph:density} even if
$w_v$ has the sign opposite to that core.

Put $S_{ij}=W_iR_2W_j^T$. Substitution of
\eqref{graph:residual-equation} in the core equations gives
\begin{equation}\label{graph:core-equation}
 (A_i+S_{ii})\mathbf w_i
 =\lambda_2\mathbf w_i-S_{ij}\mathbf w_j
 -c\delta(\mathbf x_i+W_iR_2\mathbf x_{\mathcal R}),\qquad j\ne i.
\end{equation}
The cross term has a favorable sign: it is nonnegative in the positive
core equation, and becomes nonnegative after multiplying the negative
core equation by $-1$. Multiply by a positive unit Perron vector of
$A_i+S_{ii}$. The denominator from $|\mathbf w_i|$ is at least
$\theta M$ times the sum of its coordinates. The vector in parentheses
has uniformly bounded coordinates, as
$\norm{W_iR_2}=O(k^{-1/2})$. Thus the forcing changes the resulting
quotient by at most $O(\sqrt n\,\delta)=o(1/k)$, and
\begin{equation}\label{graph:core-radius}
 \lambda_2\le\rho(A_i+S_{ii})+o(1/k).
\end{equation}
The dense core is connected, so the positive Perron vector used here
exists; adding $S_{ii}\ge0$ preserves irreducibility.

The resolvent expansion and $\norm{W_i}^2=O(k)$ give
\[
 S_{ii}=\lambda_2^{-1}W_iW_i^T+O_{\mathrm{op}}(1/k),\qquad
 \norm{S_{ii}}=O(1).
\]
Write $a_i=|C_i|$ and let $\mathbf u$ be the uniform unit vector on
$a_i$ coordinates. Since $A_i\le J-I$ entrywise, Perron monotonicity
and the spectral separation of the uniform direction yield
\begin{equation}\label{graph:rank-one}
 \rho(A_i+S_{ii})
 \le\rho(J-I+S_{ii})
 \le a_i-1+\mathbf u^TS_{ii}\mathbf u+O(1/k).
\end{equation}
Here is an explicit justification of the last error. For a unit vector
$t\mathbf u+\mathbf v$ with $\mathbf v\perp\mathbf u$, subtract
$a_i-1+\mathbf u^TS_{ii}\mathbf u$ from its quadratic form. If
$K=\norm{S_{ii}}=O(1)$, the difference is at most
$-(a_i-2K)\norm{\mathbf v}^2+2K\norm{\mathbf v}$, whose maximum is
$K^2/(a_i-2K)=O(1/k)$.
Finally,
\[
 \mathbf u^TS_{ii}\mathbf u
 =\frac{\sum_v d_i(v)^2}{a_i\lambda_2}+O(1/k)
 =\sigma_i+O(1/k),
\]
because $a_i,\lambda_2=k+O(1)$ and $\sigma_i=O(1)$.
Combine \eqref{graph:core-radius}--\eqref{graph:rank-one} for both
cores and use $\lambda_1-\lambda_2=\delta$ to obtain
\eqref{graph:pole-bound}.
\end{proof}


\section{Completing two dense terminal graphs}
\label{dense:section}

For disjoint rooted graphs $(H_1,a_1)$ and $(H_2,a_2)$, denote by
$\mathcal D(H_1,H_2;L)$ the graph obtained by joining the roots by a
path of $L$ edges, all of whose internal vertices are new.  Thus
$\mathcal D(K_k,K_k;L)=\Kite(k,L)$, and both graphs have
$2k+L-1$ vertices when $|V(H_1)|=|V(H_2)|=k$.
We prove a comparison that does not require the two rooted graphs to
be isomorphic.

\begin{proposition}\label{dense:completion}
Suppose that $k\ge1000$, $L\ge10000k$, and $H_1,H_2$ have order $k$
and minimum degree at least $k-1-k/10$.  If at least one of these
graphs is incomplete, then
\[
 \delta\bigl(\mathcal D(H_1,H_2;L)\bigr)
 >\delta\bigl(\Kite(k,L)\bigr).
\]
\end{proposition}

For a graph $H$ rooted at $a$, put $B_H=A(H-a)$ and let
$\mathbf a_H$ be its root-incidence vector: its coordinate at
$v\in V(H-a)$ equals one if $av\in E(H)$ and zero otherwise.
Define the inverse rooted response by
\begin{equation}\label{dense:response}
 F_H(z)=z-\mathbf a_H^T(zI-B_H)^{-1}\mathbf a_H,
 \qquad z>\rho(B_H).
\end{equation}
Here and below $\rho$ denotes the spectral radius of a nonnegative
symmetric matrix.  Direct differentiation gives
\begin{equation}\label{dense:response-derivatives}
 \begin{split}
 F_H'(z)&=1+\mathbf a_H^T(zI-B_H)^{-2}\mathbf a_H,\\
 F_H''(z)&=-2\mathbf a_H^T(zI-B_H)^{-3}\mathbf a_H.
 \end{split}
\end{equation}
We use the path functions $r(z)$, $s_L(z)$, and $h_L(z)$ from
Lemma~\ref{pre:path}; in particular,
\begin{equation}\label{dense:path-identity}
 s_L(z)-r(z)=-h_L(z)r(z)^L,
 \qquad
 h_L(z)=\frac{(1-r(z)^2)r(z)^{L-1}}{1-r(z)^{2L}}.
\end{equation}

\begin{lemma}\label{dense:core-estimates}
Let $H$ have order $k\ge1000$, minimum degree at least $k-1-k/10$,
and $M$ missing edges.  There is a unique number
$\alpha_H>\rho(H)$ such that $F_H(\alpha_H)=r(\alpha_H)$.
Write $\alpha_K$ for this parameter when $H=K_k$. By the rooted
clique equation in Lemma~\ref{bench:kite} and uniqueness,
$\alpha_K=\kappa_k$ and $r(\alpha_K)^{-1}=\beta_k$. Then
\begin{align}
 0<\alpha_H-\rho(H)&\le r(\alpha_H)\le\frac2k,
 &\rho(H)-\rho(B_H)&\ge\frac18,
 \label{dense:root-gap}\\
 \alpha_K-\alpha_H&\ge\frac{M}{4k},
 \label{dense:missing-deficit}\\
 F_H'(\alpha_H)&\le4k,
 &F_H'(\alpha_H)&\le k+12M+6.
 \label{dense:slope}
\end{align}
For every $z\ge\alpha_H-1/16$, the response is defined and satisfies
\begin{equation}\label{dense:log-slope}
 |F_H''(z)|\le32\bigl(F_H'(z)-1\bigr)\le32F_H'(z).
\end{equation}
\end{lemma}

\begin{proof}
We first establish estimates that remain valid as missing edges are
inserted with weights between zero and one.  Write the resulting
weighted adjacency matrix as $A=J-I-E$, where $E$ is symmetric and
nonnegative and every row sum of $E$ is at most $k/10$.
For $0\le t\le2/k$, let $\xi$ and $\boldsymbol\varphi$ be the Perron
root and positive unit Perron vector of
$A+t\mathbf e_a\mathbf e_a^T$.  These matrices are irreducible:
their underlying graphs have minimum degree greater than
$(k-1)/2$.  Row sums give
\[
 \frac9{10}k-1\le\xi\le k-1+\frac2k.
\]
Put $T=\xi+1$ and $S=\sum_j\varphi_j$.  The coordinate equations imply
\[
 \varphi_{\max}\le\frac{S}{T-t},
 \qquad
 \varphi_{\min}\ge\frac{S-(k/10)\varphi_{\max}}{T}.
\]
Consequently
\begin{equation}\label{dense:coordinates}
 \frac{\varphi_{\min}}{\varphi_{\max}}
 \ge 1-\frac{k/10+t}{T}>\frac34,
 \qquad
 \frac1{2\sqrt{k}}\le\varphi_j\le\frac2{\sqrt{k}}.
\end{equation}
The same calculation on the root-deleted graph uses
$T\ge(9/10)k-1$ and gives these latter coordinate bounds for its
unit Perron vector $\mathbf v$ as well.
During the interpolation, $B_H$ and $\mathbf a_H$ denote the current
weighted nonroot block and root-incidence vector.  The nonroot
coordinate equation yields
\[
 (\xi I-B_H)\boldsymbol\varphi_{H-a}
   =\mathbf a_H\varphi_a.
\]
Multiplication by $\mathbf v^T$, followed by Cauchy--Schwarz in the
denominator, gives
\[
 \xi-\rho(B_H)
 =\frac{(\mathbf v^T\mathbf a_H)\varphi_a}
        {\mathbf v^T\boldsymbol\varphi_{H-a}}
 \ge\frac{(9/10)k-1}{4k}>\frac18.
\]
Taking $t=0$ proves the root-deletion estimate in
\eqref{dense:root-gap}.

The Perron equation for $H$ shows that $F_H(\rho(H))=0$.
Since $F_H'>0$ and $r'<0$, the function $F_H-r$ is strictly increasing
on $[\rho(H),\infty)$, is negative at the left endpoint, and tends
to infinity.  Its unique zero is $\alpha_H$.  Eliminating the
nonroot coordinates also proves the equivalent identity
\begin{equation}\label{dense:ray-perron}
 \alpha_H
 =\rho\bigl(A(H)+r(\alpha_H)\mathbf e_a\mathbf e_a^T\bigr).
\end{equation}
Indeed the eliminated vector is positive, so its eigenvalue is the
Perron root.  The variational principle now gives
$0<\alpha_H-\rho(H)\le r(\alpha_H)$.
Since $\alpha_H\ge(9/10)k-1$, we have
$r(\alpha_H)\le2/k$.  Thus the preceding coordinate estimates apply
to the unit Perron vector in \eqref{dense:ray-perron}.

Insert all $M$ missing edges with a common weight $s\in[0,1]$,
and denote the parameter and unit Perron vector by
$\alpha_s$ and $\boldsymbol\varphi(s)$.
The eigenvalue is simple, and differentiating
\eqref{dense:ray-perron} gives
\[
 \alpha_s'
 =\frac{2\sum_{uv\text{ missing}}\varphi_u(s)\varphi_v(s)}
        {1-r'(\alpha_s)\varphi_a(s)^2}.
\]
The numerator is at least $M/(2k)$.  Moreover,
$|r'(\alpha_s)|\le4/k^2$ and $\varphi_a(s)^2\le4/k$, so the
denominator is at most two.  Integration proves
\eqref{dense:missing-deficit}.

The eigenvector in \eqref{dense:ray-perron} with root coordinate one is
\[
 \left(1,(\alpha_H I-B_H)^{-1}\mathbf a_H\right)^T.
\]
Its squared norm is $F_H'(\alpha_H)$, whence
$F_H'(\alpha_H)=1/\varphi_a^2\le4k$.
For the sharper estimate, first assume $M\le k$ and put
$\alpha=\alpha_H$.  The average degree and Perron monotonicity give
\[
 k-1-\frac{2M}{k}\le\alpha\le\alpha_K<k.
\]
The last strict inequality follows from the maximum row sum
$k-1+r(\alpha_K)<k$ in \eqref{dense:ray-perron}.
For the clique,
\[
 r(\alpha_K)^{-1}
   =\frac{k-1}{\alpha_K-k+2}<k-1.
\]
Since $r(z)^{-1}$ is increasing,
$\alpha+1-r(\alpha)=1+r(\alpha)^{-1}<k$.
The root coordinate equation and
$\varphi_{\max}/\varphi_{\min}\le2$ imply
\[
 \frac{S}{\varphi_a}
 =\alpha+1-r(\alpha)
   +\sum_{j:\,aj\text{ missing}}\frac{\varphi_j}{\varphi_a}
 \le k+2M.
\]
As $1=\sum_j\varphi_j^2\le\varphi_{\max}S$ and
$\varphi_{\max}\le S/(\alpha+1-r(\alpha))$, it follows that
\[
 F_H'(\alpha)
 \le\frac{(k+2M)^2}{k-2M/k-2/k}
 \le k+12M+6.
\]
The denominator is positive.  After clearing it, the last inequality
is equivalent to positivity of
\[
 8kM-4M^2-2M+6k-2
 -\frac{24M^2+36M+12}{k}.
\]
For $0\le M\le k$ and $k\ge1000$, this expression is at least
$(4k-27)M+6k-3>0$.  If $M\ge k$, the already established bound
$F_H'(\alpha_H)\le4k$ is stronger than the claimed estimate.
Finally,
$\alpha_H-\rho(B_H)\ge1/8$, so for
$z\ge\alpha_H-1/16$ the resolvent $(zI-B_H)^{-1}$ is positive
definite with norm at most $16$.  Applying this to
\eqref{dense:response-derivatives} proves
\eqref{dense:log-slope}.
\end{proof}

\begin{lemma}\label{dense:clique-benchmark}
For $k\ge1000$ and $L\ge k$,
\begin{equation}\label{dense:clique-parameter}
 \frac1{k^2}\le\alpha_K-(k-1)\le\frac2{k^2},
\end{equation}
and
\begin{equation}\label{dense:sharp-benchmark}
 \delta\bigl(\Kite(k,L)\bigr)
 \le\frac{2h_L(\alpha_K)}{k-2}.
\end{equation}
In addition,
\begin{equation}\label{dense:coarse-benchmark}
 \delta\bigl(\Kite(k,L)\bigr)
 \le2h_L(\alpha_K)\le2(k-2)^{-(L-1)}.
\end{equation}
\end{lemma}

\begin{proof}
Write $\alpha=\alpha_K$, $r=r(\alpha)$, $h=h_L(\alpha)$,
$\varepsilon=\alpha-(k-1)$, and
$\lambda_j^K=\lambda_j(\Kite(k,L))$.
The clique response is
\[
 F_{K_k}(z)=z-\frac{k-1}{z-k+2},
\]
so $F_{K_k}(\alpha)=r$ is equivalent to
\begin{equation}\label{dense:clique-epsilon}
 \varepsilon(k+\varepsilon)=r(1+\varepsilon).
\end{equation}
We have $0<\varepsilon\le r\le2/k$, and
$1/k<r\le1/(k-2)$ because $k-1<\alpha<k$.
Thus \eqref{dense:clique-epsilon} gives
$\varepsilon>1/k^2$ and
\[
 \varepsilon\le\frac{1+2/k}{k(k-2)}\le\frac2{k^2},
\]
proving \eqref{dense:clique-parameter}.

Label the joining path $a_0,\ldots,a_L$.
Give a vector $\mathbf u$ the coordinates $u_{a_j}=r^j$,
and give each nonroot clique vertex its own root coordinate divided
by $\alpha-k+2$.  Let $\mathbf v$ be the reflected vector.
Substitution into the adjacency equations gives
\[
 (A-\alpha I)\mathbf u=e\mathbf e_{a_L},
 \qquad
 (A-\alpha I)\mathbf v=e\mathbf e_{a_0},
 \qquad e=(1-r^2)r^{L-1}.
\]
The positive-vector bound applied to $\mathbf u+\mathbf v$ gives
\[
 \alpha<\lambda_1^K\le\alpha+\frac{e}{1+r^L}.
\]
Reflection fixes the positive Perron vector, so
$\mathbf u-\mathbf v$ is orthogonal to it.  Its Rayleigh quotient,
using $\lVert\mathbf u-\mathbf v\rVert_2^2\ge2(1-r^L)^2$, gives
\[
 \lambda_2^K
 \ge\alpha-\frac{2e(1-r^L)}{\lVert\mathbf u-\mathbf v\rVert_2^2}
 \ge\alpha-\frac{e}{1-r^L}.
\]
Consequently
\[
 \delta\bigl(\Kite(k,L)\bigr)\le2h,
 \qquad
 \lambda_1^K,\lambda_2^K\in[\alpha-2h,\alpha+2h].
\]
Since $h\le r^{L-1}\le(k-2)^{-(L-1)}$, this proves
\eqref{dense:coarse-benchmark}.

We shall use the uniform estimate
\begin{equation}\label{dense:tiny}
 L(k-2)^{-(L-1)}
 \le k(k-2)^{-(k-1)}\le k^{-20}
 \qquad(L\ge k\ge1000).
\end{equation}
The expression decreases with integral $L\ge k$, and
$k-2\ge\sqrt{k}$ proves the second inequality.
It follows from \eqref{dense:clique-parameter} that
$\lambda_2^K>k-1$ and that
$z-k+2\le1+3/k^2$ on $[\alpha-2h,\alpha+2h]$.
For $Q=F_{K_k}-r$, we therefore have
\begin{equation}\label{dense:clique-Q-slope}
 Q'(z)\ge F_{K_k}'(z)
   =1+\frac{k-1}{(z-k+2)^2}\ge k-1
\end{equation}
throughout this interval.

After deleting both roots, the remaining components are two
$(k-1)$-cliques and the internal path.  Their spectral radii are
below $k-1<\lambda_2^K$.  The two-root Schur complement and reflection
symmetry therefore identify the two eigenvalues by
\begin{align*}
 Q(\lambda_1^K)
 &=h_L(\lambda_1^K)\bigl(1-r(\lambda_1^K)^L\bigr),\\
 Q(\lambda_2^K)
 &=-h_L(\lambda_2^K)\bigl(1+r(\lambda_2^K)^L\bigr).
\end{align*}
Indeed the symmetric root vector belongs to the larger eigenvalue;
the Schur matrix has inertia $(\text{negative},0)$ at the second,
so its null vector there is antisymmetric.
For $z\ge k-1$, differentiation of \eqref{dense:path-identity} gives
\begin{equation}\label{dense:log-h}
 0<-\frac{d}{dz}\log h_L(z)
 =\frac{(L-1)-2r(z)^2/(1-r(z)^2)
             +2Lr(z)^{2L}/(1-r(z)^{2L})}
        {\sqrt{z^2-4}}
 \le\frac{4L}{k}.
\end{equation}
In particular, \eqref{dense:tiny} implies
\[
 \frac{h_L(\lambda_2^K)}{h_L(\alpha)}
 \le\exp(8Lh/k)\le1+\frac1{50k},
 \qquad
 r(\lambda_2^K)^L\le(k-2)^{-L}\le\frac1{100k}.
\]
Integrating \eqref{dense:clique-Q-slope} between the two eigenvalues,
and using that $h_L$ decreases, yields
\[
 \delta\bigl(\Kite(k,L)\bigr)
 \le\frac{2h_L(\lambda_2^K)(1+r(\lambda_2^K)^L)}{k-1}
 \le\frac{2h_L(\alpha_K)}{k-2}.
\]
This proves \eqref{dense:sharp-benchmark}.
\end{proof}

\begin{proof}[Proof of Proposition~\ref{dense:completion}]
Let $G=\mathcal D(H_1,H_2;L)$ and suppose, towards a contradiction,
that $\delta=\delta(G)\le\delta(\Kite(k,L))$.
Write $\alpha_i=\alpha_{H_i}$, let $M_i$ be the number of missing
edges in $H_i$, and put $M=\max(M_1,M_2)\ge1$.
By \eqref{dense:coarse-benchmark} and \eqref{dense:tiny},
\begin{equation}\label{dense:small-gap}
 \delta\le2(k-2)^{-(L-1)},
 \qquad L\delta\le2k^{-20}.
\end{equation}
In this proof $\lambda_1,\lambda_2$ are the two largest eigenvalues
of $G$.

Since $\lambda_1\ge\rho(H_i)$, the root-deletion estimate gives
$\lambda_2>\rho(H_i-a_i)$ and $\lambda_2>2$.
Deleting both roots thus leaves a matrix with spectral radius below
$\lambda_2$.  Two-vertex interlacing implies
$\lambda_3<\lambda_2$, so the second eigenvalue is simple.
On $[\lambda_2,\lambda_1]$, elimination of all nonroot vertices gives
\begin{equation}\label{dense:schur}
 \mathcal F(z)=
 \begin{pmatrix}
 F_{H_1}(z)-s_L(z)&-h_L(z)\\
 -h_L(z)&F_{H_2}(z)-s_L(z)
 \end{pmatrix}.
\end{equation}
Its inertia is $(\text{negative},0)$ at $\lambda_2$ and
$(0,\text{positive})$ at $\lambda_1$.
Furthermore $\mathcal F'(z)\succeq I$.
The path contribution to this derivative has norm at most
$(\lambda_2-2)^{-2}$: the attachment columns at the two distinct
ends of the internal path are orthonormal, and the path resolvent
has norm at most $(\lambda_2-2)^{-1}$.

If $z\ge\alpha_i$, then $F_{H_i}'(z)\le4k$ by
\eqref{dense:response-derivatives} and \eqref{dense:slope}.
If $z<\alpha_i$, we still have
\[
 z\ge\lambda_2\ge\rho(H_i)-\delta
   \ge\alpha_i-2/k-\delta.
\]
This lies above $\alpha_i-1/16$.
Integrating \eqref{dense:log-slope} gives
$F_{H_i}'(z)\le4k\exp(32(2/k+\delta))<8k$.
Including the path contribution, we obtain
\begin{equation}\label{dense:rough-matrix-slope}
 I\preceq\mathcal F'(z)\preceq10kI
 \qquad(\lambda_2\le z\le\lambda_1).
\end{equation}
Lemma~\ref{pre:schur} therefore gives
\[
 0\le\mathcal F_{ii}(\lambda_1)\le10k\delta,
 \qquad h_L(\lambda_1)\le5k\delta.
\]
Using \eqref{dense:path-identity}, we find
\[
 -h_L(\lambda_1)r(\lambda_1)^L
 \le F_{H_i}(\lambda_1)-r(\lambda_1)\le10k\delta.
\]
The function $F_{H_i}-r$ has derivative at least one and zero
$\alpha_i$.  Hence
\begin{equation}\label{dense:resonance}
 |\lambda_1-\alpha_i|\le10k\delta\qquad(i=1,2).
\end{equation}
Thus proximity of the two rooted parameters follows from the
hypothetical small gap; it has not been assumed.

Every $z\in[\lambda_2,\lambda_1]$ is within $11k\delta$ of each
$\alpha_i$.  Since $M_i\le k^2/20$,
$k+12M_i+6\le k^2$.  Equations \eqref{dense:slope} and
\eqref{dense:log-slope} show that the increase in either scalar
derivative from its value at $\alpha_i$ is at most
\[
 k^2\bigl(\exp(352k\delta)-1\bigr)<1,
\]
where the last inequality follows from \eqref{dense:small-gap}.
The path contribution is also less than one.  It follows that
\begin{equation}\label{dense:sharp-matrix-slope}
 \mathcal F'(z)\preceq\bigl(k+20(M+1)\bigr)I.
\end{equation}
A second application of Lemma~\ref{pre:schur} gives
\begin{equation}\label{dense:gap-lower}
 \delta\ge\frac{2h_L(\lambda_1)}{k+20(M+1)}.
\end{equation}

Apply \eqref{dense:missing-deficit} to a core with $M$ missing edges.
Together with \eqref{dense:resonance}, this gives
\begin{equation}\label{dense:radius-loss}
 \lambda_1\le\alpha_K-\frac{M}{4k}+10k\delta<\alpha_K,
\end{equation}
since $10k\delta<1/(8k)\le M/(8k)$.
For $z\in[\lambda_1,\alpha_K]$,
\[
 -\frac{d}{dz}\log r(z)=\frac1{\sqrt{z^2-4}}\ge\frac1k,
\]
as $\alpha_K<k$.  Also $r(\lambda_1)\le2/k$.
In the ratio $h_L(\lambda_1)/h_L(\alpha_K)$, the factor involving
$1-r^{2L}$ is at least one, while the logarithm of the ratio of the
$1-r^2$ factors is at least $-20/k^2$.  Therefore
\begin{align*}
 \log\frac{h_L(\lambda_1)}{h_L(\alpha_K)}
 &\ge\frac{(L-1)(\alpha_K-\lambda_1)}{k}-\frac{20}{k^2}\\
 &\ge\frac{M(L-1)}{4k^2}-10L\delta-\frac{20}{k^2}.
\end{align*}
Combining \eqref{dense:sharp-benchmark} and
\eqref{dense:gap-lower}, we conclude that
\begin{align*}
 \log\frac{\delta(G)}{\delta(\Kite(k,L))}
 &\ge\frac{M(L-1)}{4k^2}-10L\delta-\frac{20}{k^2}
       -\log\frac{k+20(M+1)}{k-2}\\
 &\ge\frac{M(L-1)}{4k^2}-10L\delta
       -\frac{20(M+1)+3}{k}-\frac{20}{k^2}.
\end{align*}
Here we used
\[
 \log\frac{k+20(M+1)}{k-2}
 =\log\left(1+\frac{20(M+1)}k\right)
  -\log\left(1-\frac2k\right)
 \le\frac{20(M+1)+3}{k}.
\]
Finally, $L-1\ge9999k$ and $L\delta\le2k^{-20}$ show that the
last lower bound is at least
\[
 \frac{9919M-92}{4k}-\frac{20}{k^2}-20k^{-20}
 >\frac{2400M}{k}>0.
\]
This contradicts $\delta(G)\le\delta(\Kite(k,L))$ and proves the
proposition.
\end{proof}

\section{Proof of the graph theorem}\label{sec:graph-conclusion}

\begin{proof}[Proof of Theorem~\ref{thm:graphs}]
Suppose there are non-double-kite minimizers at arbitrarily large
orders, and choose a sequence of them. Apply the preceding graph
reductions. With their notation, the graph
$K=\Kite(k,L+b+q)$ has the same order as $G$. The sharp clique
benchmark in Lemma~\ref{dense:clique-benchmark} gives
\[
 \delta(K)\le\frac2k(1+o(1))\beta_k^{-(L+b+q-1)},
 \qquad\beta_k=k-1+O(1/k).
\]
Dividing the transmission bound \eqref{graph:transmission-bound} by
this estimate, and taking logarithms, gives
\begin{align}
 \log\frac{\delta(G)}{\delta(K)}
 \ge{}&\log\theta^2+\log(d_1d_2)+(b+q)\log\beta_k\notag\\
 &-(L-1)\log(\lambda_1/\beta_k)+o(1).
 \label{graph:comparison-exact}
\end{align}
The spectral bound \eqref{graph:pole-bound} implies
\[
 \log(\lambda_1/\beta_k)
 \le\frac{\min(\sigma_1,q+\sigma_2)}k+O(1/k^2).
\]
Since $L/k=2\log k+o(\log k)$, and $b+q=O(1)$, it follows that
\begin{align}
 \log\frac{\delta(G)}{\delta(K)}
 \ge{}&\log\theta^2+\log(d_1d_2)+(b+q)\log k\notag\\
 &-(2\log k+o(\log k))\min(\sigma_1,q+\sigma_2)+o(1).
 \label{graph:comparison}
\end{align}
In particular the $O(1/k)$ error in the spectral estimate contributes
only $O(\log k/k)=o(1)$ here.

Pass to a subsequence on which the bounded integers $b,q$ are fixed.
Label the $b$ residual vertices outside the shortest core path in any
order. Passing to a further subsequence, all their attachment densities,
the two endpoint densities $d_i/k$, and the two numbers
$\log d_i/\log k$ converge. Write
\[
 d_i/k\to t_i,\qquad \frac{\log d_i}{\log k}\to\alpha_i\in[0,1],
 \qquad \sigma_i\to t_i^2+H_i,
\]
where $H_i$ is the sum of squared limiting densities of off-path
vertices attached to core $i$. If $t_i>0$, then $\alpha_i=1$.
By \eqref{graph:density}, each limiting density is at most $\theta$.
Each off-path vertex meets at most one core, so
$H_1+H_2\le b\theta^2$.

The limiting coefficient of $\log k$ on the right of
\eqref{graph:comparison} is
\[
 P=b+q+\alpha_1+\alpha_2
 -2\min(t_1^2+H_1,q+t_2^2+H_2).
\]
Using $2\min(u,v)\le u+v$ gives the useful decomposition
\begin{equation}\label{graph:coefficient}
 P\ge b(1-\theta^2)+
 (\alpha_1-t_1^2)+(\alpha_2-t_2^2).
\end{equation}
Each summand on the right is nonnegative. If $t_i=0$, its contribution
is $\alpha_i$; if $t_i>0$, it is $1-t_i^2\ge1-\theta^2>0$.
Consequently $P>0$ if $b>0$, if either $t_i>0$, or if either
$\alpha_i>0$. In any of these cases \eqref{graph:comparison} tends
to $+\infty$, contradicting minimality, which requires the logarithmic
ratio to be nonpositive. If all these quantities vanish but $q>0$,
then $H_1=H_2=0$ and the exact expression gives $P=q>0$, again a
contradiction.

The only possible remaining subsequence therefore has $b=q=0$ and
$\alpha_1=\alpha_2=0$. Now $d_i=k^{o(1)}$. There are no off-path
vertices, so
$\sigma_i=d_i^2/k^2=k^{-2+o(1)}$; its entire logarithmic penalty in
\eqref{graph:comparison} is $o(1)$. We conclude that
\[
 \frac{\delta(G)}{\delta(K)}
 \ge\theta^2d_1d_2(1-o(1)).
\]
If the integer product $d_1d_2$ is at least two on a further
subsequence, this exceeds one because $2\theta^2>1$. Thus
$d_1=d_2=1$ eventually.

Every residual vertex now lies on the shortest path between the two
cores. That path is induced; shortestness excludes every other
core-to-path edge. The two endpoint attachments are single edges, the
cores have equal order $k$, and there are no edges between the cores.
Thus $G=\mathcal D(G[C_1],G[C_2];L)$ in the notation of
Section~\ref{dense:section}. Lemma~\ref{graph:cores} gives the
required core minimum degrees. Since $L/(k\log k)\to2$, eventually
$k\ge1000$ and $L\ge10000k$. Proposition~\ref{dense:completion}
forces both cores to be complete, contradicting the choice of the
sequence.

The same reasoning applies to any sequence of exceptional minimizers.
Their existence at unbounded orders is therefore impossible, proving
the existence of one $N_{\mathcal G}$ valid for every minimizer.
The parameter asymptotics follow from $k=k_0+O(1)$,
\eqref{graph:scale-eq}, and $n=2k+L-1$.
\end{proof}

\section{Scope of the conclusions}
Theorems~\ref{thm:trees} and \ref{thm:graphs} identify every minimizing
graph in their respective classes at all sufficiently large orders.
The proofs give existence of thresholds and do not determine the
remaining finite ranges. They also leave open whether the optimizing
integer parameter is unique at every sufficiently large order.
All comparisons above are between graphs of the same order; none of
the structural conclusions requires a numerical enumeration.

\bibliographystyle{plainnat}
\bibliography{spectral_gap_trees_and_graphs}
\end{document}
